\documentclass[t]{beamer}

\subtitle{Section 2.1: Binary Operations}

\usepackage{amsthm,amsmath,amsfonts,hyperref,graphicx,color,multicol,soul}
\usepackage{enumitem,tikz,tikz-cd,setspace,mathtools}

%%%%%%%%%%
%Beamer Template Customization
%%%%%%%%%%
\setbeamertemplate{navigation symbols}{}
\setbeamertemplate{theorems}[ams style]
\setbeamertemplate{blocks}[rounded]

\definecolor{Blu}{RGB}{43,62,133} % UWEC Blue
\setbeamercolor{structure}{fg=Blu} % Titles

%Unnumbered footnotes:
\newcommand{\blfootnote}[1]{%
	\begingroup
	\renewcommand\thefootnote{}\footnote{#1}%
	\addtocounter{footnote}{-1}%
	\endgroup
}

%%%%%%%%%%
%TikZ Stuff
%%%%%%%%%%
\usetikzlibrary{arrows}
\usetikzlibrary{shapes.geometric}
\tikzset{
	smaller/.style={
		draw,
		regular polygon,
		regular polygon sides=3,
		fill=white,
		node distance=2cm,
		minimum height=1in,
		line width = 2pt
	}
}
\tikzset{
	smsquare/.style={
		draw,
		regular polygon,
		regular polygon sides=4,
		fill=white,
		node distance=2cm,
		minimum height=1in,
		line width = 2pt
	}
}


%%%%%%%%%%
%Custom Commands
%%%%%%%%%%

\newcommand{\C}{\mathbb{C}}
\newcommand{\quats}{\mathbb{H}}
\newcommand{\N}{\mathbb{N}}
\newcommand{\Q}{\mathbb{Q}}
\newcommand{\R}{\mathbb{R}}
\newcommand{\Z}{\mathbb{Z}}

\newcommand{\ds}{\displaystyle}

\newcommand{\fn}{\insertframenumber}

\newcommand{\id}{\operatorname{id}}
\newcommand{\im}{\operatorname{im}}
\newcommand{\Aut}{\operatorname{Aut}}
\newcommand{\Inn}{\operatorname{Inn}}

\newcommand{\blank}[1]{\underline{\hspace*{#1}}}

\newcommand{\abar}{\overline{a}}
\newcommand{\bbar}{\overline{b}}
\newcommand{\cbar}{\overline{c}}

\newcommand{\nml}{\unlhd}

%%%%%%%%%%
%Custom Theorem Environments
%%%%%%%%%%
\theoremstyle{definition}
\newtheorem{exercise}{Exercise}
\newtheorem{question}[exercise]{Question}
\newtheorem{warmup}{Warm-Up}
\newtheorem*{defn}{Definition}
\newtheorem*{exa}{Example}
\newtheorem*{disc}{Group Discussion}
\newtheorem*{nb}{Note}
\newtheorem*{recall}{Recall}
\renewcommand{\emph}[1]{{\color{blue}\texttt{#1}}}

\definecolor{Gold}{RGB}{237, 172, 26}
%Statement block
\newenvironment{statementblock}[1]{%
	\setbeamercolor{block body}{bg=Gold!20}
	\setbeamercolor{block title}{bg=Gold}
	\begin{block}{\textbf{#1.}}}{\end{block}}
\newenvironment{thm}[1]{%
	\setbeamercolor{block body}{bg=Gold!20}
	\setbeamercolor{block title}{bg=Gold}
	\begin{block}{\textbf{Theorem #1.}}}{\end{block}}


%%%%%%%%%%
%Custom Environment Wrappers
%%%%%%%%%%
\newcommand{\enumarabic}[1]{
	\begin{enumerate}[label=\textbf{\arabic*.}]
		#1
	\end{enumerate}
}
\newcommand{\enumalph}[1]{
	\begin{enumerate}[label=(\alph*)]
		#1
	\end{enumerate}
}
\newcommand{\bulletize}[1]{
	\begin{itemize}[label=$\bullet$]
		#1
	\end{itemize}
}
\newcommand{\circtize}[1]{
	\begin{itemize}[label=$\circ$]
		#1
	\end{itemize}
}
\newcommand{\slide}[1]{
	\begin{frame}{\fn}
		#1
	\end{frame}
}
\newcommand{\slidec}[1]{
\begin{frame}[c]{\fn}
	#1
\end{frame}
}
\newcommand{\slidet}[2]{
	\begin{frame}{\fn\ - #1}
		#2
	\end{frame}
}


\newcommand{\startdoc}{
		\title{Math 425: Abstract Algebra 1}
		\author{Mckenzie West}
		\date{Last Updated: \today}
		\begin{frame}
			\maketitle
		\end{frame}
}

\newcommand{\topics}[2]{
	\begin{frame}{\insertframenumber}
		\begin{block}{\textbf{Last Section.}}
			\begin{itemize}[label=--]
				#1
			\end{itemize}
		\end{block}
		\begin{block}{\textbf{This Section.}}
			\begin{itemize}[label=--]
				#2
			\end{itemize}
		\end{block}
	\end{frame}
}

\begin{document} 
	\startdoc

	\topics{ 
		\item Disjoint Cycles
		\item Transpositions
		\item Even vs Odd Permutations
	}{
		\item Binary Operations
		\item Monoids
	}

\slide{
	\begin{defn}
		A \emph{binary operation}, $*$ on a set $S$ is a function that associates to each ordered pair $(a,b)\in S\times S$ an element of $S$ which we call $a*b$.
		
		Since we know that $a*b\in S$ for all $a,b\in S$, we say that the binary operation is \emph{closed under $*$}.
	\end{defn}\vskip 1in
%	\begin{exa}\begin{itemize}[label=$\bullet$]
%			\item $+,-,\cdot,\div$ are all binary operations on $\R$
%			\item $+,-,\cdot$ are all binary operations on $\Z$ and on $\Z_n$
%			\item $\circ$ is a binary operation on $S_n$
%		\end{itemize}
%	\end{exa}
	\begin{nb}
		 Sometimes we write $(S,*)$ to mean $*$ is a binary operation on $S$.
	\end{nb}
}

\slide{
	\begin{defn}
		A binary operation $*$ on $S$ is \emph{associative} if 
		\[a*(b*c)=(a*b)*c,\]
		for all $a,b,c\in S$.
	\end{defn}
}
\slide{
	\begin{defn}	
		A binary operation $*$ on $S$ is \emph{commutative} if 
		\[a*b=b*a,\]
		for all $a,b\in S$.
	\end{defn}
}
\slide{
	\begin{defn}
		An element $e\in S$ is called an \emph{identity} (or \emph{unity}) for the binary operation $*$ if 	
			\[a*e=e*a=a,\]
		for all $a\in S$.
	\end{defn}
}
\slide{
	\begin{statementblock}{Theorem 2.1.1}
		If a binary operation $*$ on a set $S$ has an identity, then it is unique.
	\end{statementblock}
}

\slide{
	\begin{defn}
		A set $S$ along with a binary operation $*$ is called an \emph{monoid} if $*$ is associative and has an identity.
		
		\vskip 1em
		
		If $(S,*)$ is also commutative, then we say $S$ is a \emph{commutative monoid}.
	\end{defn}
	\begin{exa}
		\bulletize{
			\item $(\Z,\cdot)$
			\item $(S_n,\circ)$
		}
	\end{exa}
}
\slide{
	\huge\textbf{Let's play}
	\begin{center}
		\includegraphics[width=3.5in]{monoid_game_level}
	\end{center}
}
\begin{frame}[fragile]
	\frametitle{\fn}
\begin{exercise}
	Let $*$ be the binary operation on $\mathbb{N}$ given by \[a*b:=a^b.\]  
	\enumalph{
		\item Associative?
		\vskip 2em
		\item Identity?
		\vskip 2em
		\item Commutative?
		\vskip 2em
	}
	Therefore, $(\mathbb{N},$\verb|^|$)$ is ...
\end{exercise}
\begin{picture}(0,0)
\put(150,140){\includegraphics[width=2.5in]{monoid_game}}
\end{picture}
\end{frame}
\slide{
\begin{exercise}
	Let $*$ be the binary operation on $\mathbb{Z}$ given by \[a*b:=a-b.\]  
	\enumalph{
		\item Associative?
		\vskip 2em
		\item Identity?
		\vskip 2em
		\item Commutative?
		\vskip 2em
	}
	Therefore, $(\mathbb{Z},-)$ is ...
\end{exercise}
\begin{picture}(0,0)
\put(150,140){\includegraphics[width=2.5in]{monoid_game}}
\end{picture}
}
\slide{
\begin{exercise}
	Let $*$ be the binary operation on $GL_2(\R)$given by \[A*B:=AB.\]  
	\enumalph{
		\item Associative?
		\vskip 2em
		\item Identity?
		\vskip 2em
		\item Commutative?
		\vskip 2em
	}
	Therefore, $(GL_2(\R),\cdot)$ is ...
\end{exercise}
\begin{picture}(0,0)
\put(150,140){\includegraphics[width=2.5in]{monoid_game}}
\put(160,0){\rotatebox{15}{$GL_2(\R)$ - $2\times 2$ invertible matrices}}
\end{picture}
}
\slide{
	\begin{defn}
		Let $(M,*)$ be a monoid.  
		
		If $x\in M$, we call $y\in M$ an \emph{inverse of $x$} if 	
			\[xy=e=yx.\]
		An element that has an inverse is called a \emph{unit}.
	\end{defn}
}
\slide{\vskip -.25in
	\begin{exercise}
		Fill in the table
		
			\def\arraystretch{1.5}
			$\begin{array}{c|c|c}
			\text{Monoid}&\quad e\quad &\quad \text{Set of Units}\quad\\
			\hline\hline
			(\Z,+)&&\\\hline
			(\R,+)&&\\\hline
			(\Z_n,+)&&\\\hline
			(M_n(\R),+)&&\\\hline
			(M_n(\R),\cdot)&&\\\hline
			(\R,\cdot)&&\\\hline
			(\Z_4,\cdot)&&\\\hline
			(\Z_n,\cdot)&&\\\hline
			(\Z_p,\cdot)&&\\\hline
			(S_3,\circ)&&\\\hline
			\end{array}$
			
		\begin{picture}(0,0)
		\put(195,132){$\leftarrow$ Matrices with real entries}
		\put(195,30){$\leftarrow$ $p$ is prime}
		\end{picture}
	\end{exercise}
}

\end{document}

